% =====================================================================
% 2. Arquitectura y decisiones tecnológicas
% =====================================================================
\section{Arquitectura y decisiones tecnológicas}
\label{sec:arquitectura}

\subsection{Decisiones}

\begin{table}[H]
\centering\small
\caption{Decisiones tecnológicas del proyecto.}
\label{tab:decisiones}
\begin{tabularx}{\textwidth}{@{}p{2.6cm}p{3.6cm}L@{}}
\toprule
\textbf{Decisión} & \textbf{Elección} & \textbf{Justificación} \\
\midrule
Motor NoSQL & MongoDB 8 (modelo documental) &
  Guarda geometrías en GeoJSON de forma nativa, tiene índices \campo{2dsphere}
  y operadores \campo{\$geoWithin} y \campo{\$geoIntersects}~\cite{mongodb}.
  Su pipeline de agregación (\campo{\$group}, \campo{\$lookup},
  \campo{\$merge}) calcula los totales por municipio dentro de la base, y
  puede escalar horizontalmente con \emph{sharding}. \\
Fuente de edificios 1 & Google Open Buildings v3 &
  Cubre América Latina con imágenes de alta resolución (inferencia de mayo
  de 2023). Trae \campo{area\_in\_meters} y \campo{confidence} (0,65 a 1,0),
  lo que permite validar las áreas y medir la sensibilidad. Licencia
  CC BY 4.0 / ODbL. \\
Fuente de edificios 2 & Microsoft Building Footprints &
  Imágenes de 2014 a 2021, licencia ODbL. Particionado por país y
  \emph{quadkey}, así que se descarga solo lo necesario. Solo trae la
  geometría: el área se calcula. \\
Límites municipales & MGN 2025 del DANE, nivel municipio &
  Fuente oficial que exige el enunciado; se filtra con los 170 códigos del
  listado PDET. \\
Sistemas de referencia & WGS84 (EPSG:4326) para almacenar; MAGNA-SIRGAS
  Origen Nacional (EPSG:9377) para áreas &
  GeoJSON y \campo{2dsphere} exigen \campo{[longitud, latitud]} en WGS84. Las
  áreas se calculan en la proyección plana oficial de Colombia, en metros. \\
Herramientas & Python, GeoPandas, Shapely, PyMongo, DuckDB, Docker Compose,
  GitHub &
  GeoPandas valida y reproyecta; DuckDB lee y filtra los archivos grandes sin
  cargarlos completos; PyMongo carga por lotes; Docker Compose levanta la
  misma base en cualquier máquina. \\
\bottomrule
\end{tabularx}
\end{table}

\subsection{Flujo de datos}

El flujo tiene cuatro etapas encadenadas por un \campo{Makefile}
(figura~\ref{fig:flujo}). Cada etapa es una carpeta de \campo{src/} con un
archivo por tarea, lee sus parámetros de \campo{config.yaml} y deja un
registro en la colección \campo{bitacora\_cargas}. Las etapas se comunican por
archivos intermedios en \campo{data/processed/} y por las colecciones de
MongoDB, así que cada una se puede volver a correr por separado.

\begin{figure}[H]
\centering
\begin{tikzpicture}[
  font=\small,
  caja/.style={draw=acento, rounded corners=3pt, fill=acento!6, align=center,
               minimum height=1.25cm, text width=3.0cm},
  db/.style={draw=gris, rounded corners=3pt, fill=gray!8, align=center, font=\footnotesize,
             minimum height=0.9cm, text width=3.2cm},
  flecha/.style={-{Stealth[length=2.2mm]}, thick, gris},
  node distance=0.55cm and 0.55cm]
  \node[caja] (m) {\textbf{make municipios}\\ MGN 2025 + PDET\\ validación, carga, índice};
  \node[caja, right=of m] (e) {\textbf{make edificios}\\ Google y Microsoft\\ prefiltro, limpieza, carga};
  \node[caja, right=of e] (a) {\textbf{make analisis}\\ asignación, agregación,\\ sensibilidad};
  \node[caja, right=of a] (s) {\textbf{Salidas}\\ CSV, tablas, mapas,\\ ranking, potencial};
  \node[db, below=of m] (dm) {\campo{municipios}\\ \campo{fuentes}};
  \node[db, below=of e] (de) {\campo{edificios\_google}\\ \campo{edificios\_microsoft}};
  \node[db, below=of a] (da) {\campo{resultados}};
  \node[db, below=of s] (ds) {\ruta{data/processed/}\\ informe};
  \draw[flecha] (m) -- (e); \draw[flecha] (e) -- (a); \draw[flecha] (a) -- (s);
  \draw[flecha] (m) -- (dm); \draw[flecha] (e) -- (de); \draw[flecha] (a) -- (da);
  \draw[flecha] (s) -- (ds);
  \draw[flecha, dashed] (dm.east) -- (de.west);
  \draw[flecha, dashed] (de.east) -- (da.west);
  \draw[flecha, dashed] (da.east) -- (ds.west);
\end{tikzpicture}
\caption{Etapas del flujo (\campo{make all = make up municipios edificios
analisis}) y colecciones que produce cada una. Todas registran su ejecución en
\campo{bitacora\_cargas}.}
\label{fig:flujo}
\end{figure}

\subsection{Organización del repositorio y del trabajo}

\begin{lstlisting}
src/common/      conexión a MongoDB y lectura de config.yaml
src/municipios/  límites PDET: descarga, cruce, geometría, esquema, carga, índices
src/edificios/   huellas: descarga y limpieza por fuente, esquema, carga, índices
src/analisis/    asignación, agregación, sensibilidad, figuras, mapas, ranking
config.yaml      parámetros (CRS, umbrales, rutas, colecciones)
data/            datos crudos y procesados (solo se versionan las salidas pequeñas)
Entrega N/       informe de cada entrega; Informe Final/ este documento
\end{lstlisting}

El equipo trabajó con una rama y un \emph{pull request} por persona y por
tarea, repartidas por archivo para que las cuatro personas avanzaran en
paralelo. La tabla~\ref{tab:reparto} resume quién hizo qué.

\begin{table}[H]
\centering\small
\caption{Reparto del trabajo por semana.}
\label{tab:reparto}
\begin{tabularx}{\textwidth}{@{}p{2.4cm}LLLL@{}}
\toprule
\textbf{Semana} & \textbf{Jorge} & \textbf{Maria Camila} & \textbf{Samuel} & \textbf{Jose Eduardo} \\
\midrule
2. Límites PDET & Descarga del MGN, cruce y verificación & Geometría y exportación & Esquema, carga, índices y evidencia & Mapa de control y tamaño BSON \\
3. Edificios & Descarga y limpieza de Google & Descarga y limpieza de Microsoft & Esquema, carga, índices y \campo{make edificios} & EDA \\
4. Análisis & Asignación, \campo{explain} y validación & Mapas, ranking, \campo{run.py} e integración & Agregación, sensibilidad y exportación & Estilo de figuras y lectura de resultados \\
\bottomrule
\end{tabularx}
\end{table}
